\section{Introduction}
Large Language Models (LLMs)~\citep{achiam2023gpt,team2023gemini,liu2024deepseek,yang2025qwen3} have demonstrated impressive reasoning abilities across diverse domains (e.g., question answering, code generation), but many practical applications involve complex scenarios where multiple agents need to interact and coordinate. For example, tasks like complex information retrieval~\citep{chang2025main,zhang2025agentorchestra}, agentic software engineering~\citep{hong2024metagpt,qian2024chatdev}, and open-ended device control~\citep{wang2024mobile,tan2024cradle} involve multiple participants working together over extended horizons. Organizing LLMs into a Multi-Agent System (MAS), where each agent specializes in a subtask or role, has become a trend to handle complex real-world tasks more effectively~\citep{tran2025multi,zhang2025landscape}.


Reinforcement Learning (RL)~\citep{sutton2018reinforcement} now plays a foundational role in LLM post-training. Despite its growing importance, the extension of RL training to multi-agent LLM systems remains largely underexplored from both \textit{algorithmic} and \textit{system} perspectives.
On the algorithmic side, while group-based RL methods like Group Relative Policy Optimization (GRPO)~\citep{shao2024deepseekmath} excel in single-agent scenarios, adapting them to multi-agent settings introduces significant challenges due to the frequent \textit{instability} observed across various scenarios~\citep{chen2025heterogeneous,zhao2025stronger,yuan2025marshal}. Agents are often invoked at different frequencies, leading to heterogeneous data distributions that greatly complicate end-to-end optimization~\citep{hong2025multi}.
On the system side, recent large-scale RL post-training frameworks, like veRL~\citep{sheng2024hybridflow}, ROLL~\citep{wang2025reinforcement}, and AReaL~\citep{fu2025areal}, provide flexible, high-throughput training pipelines for LLMs, but are largely designed around optimizing a single LLM actor. They generally lack the native support for efficient multi-agent orchestration and multiple LLMs' co-training, restricting the flexible scheduling and resource sharing required for heterogeneous agent configurations.


In this work, we theoretically identify that applying vanilla GRPO to train multi-agent LLM systems introduces systematic gradient variance and destabilizes training. We provide rigorous mathematical and empirical analysis, demonstrating that using a global advantage baseline across agents can inflate the second moment of their gradients, leading to gradient-norm explosion. 
Building on this analysis, we propose \methodname{}, a simple and stable RL training recipe for multi-agent LLM systems. \methodname{} adopts a straightforward yet effective remedy: each agent normalizes rewards using its own mean and variance. Concretely, we group action experience by agent so that each policy’s advantage estimates are normalized with respect to its own data distribution.
This calibration balances per-agent gradients, thus resulting in a dramatic reduction in variance for the policy gradient estimator. Beyond the algorithm itself, \methodname{} also provides an end-to-end RL training framework tailored for multi-agent LLM systems. It supports scalable multi-agent orchestration, flexible agent-model assignment with optional LLM sharing (e.g., co-training 7B and 3B models), per-agent configuration of optimization, and shared resource pooling for efficient scheduling of LLM actor backends. The result is a unified system that maintains well-conditioned gradients and high hardware efficiency while enabling stable co-training across multiple LLM agents.

We evaluate \methodname{} on role-specialized multi-agent systems for math reasoning and multi-turn search, using Qwen2.5~\citep{bai2025qwen2} and Qwen3~\citep{yang2025qwen3} series models, under both LLM-sharing and non-sharing settings. Across tasks and settings, \methodname{} consistently improves the performance over vanilla GRPO (e.g., +5.6\% avg@16 and +4.6\% pass@16 on math, and +15.2\% avg@16 and +13.1\% pass@16 on search). We also observe markedly improved stability, with gradient-norm spikes largely eliminated. Furthermore, \methodname{} remains highly effective under heterogeneous agent-model assignments, enabling smaller models for lower-level agents' decisions while improving overall system efficiency. 
% Our code is available at~\url{https://anonymous.4open.science/r/DrMAS-4A1E}.



